\subsection{Ordered Sets, Well-Orders, and Lattices}

\begin{definition} \label{def:ordered-set-vocabulary} \hlt{Ordered Set Vocabulary}\\
Let $S$ be a set, let $E,C\subseteq S$, and let $\leq_S$ be a relation on $S$. Write $x<_Sy$ iff $x\leq_Sy$ and $x\neq y$.
\begin{enumerate}[label=(\roman*)]
\setlength{\itemsep}{0pt}
\item $\leq_S$ is a \hlt{partial order} if it is reflexive, antisymmetric, and transitive; then $(S,\leq_S)$ is a \hlt{poset}.
\item $\leq_S$ is an \hlt{order} if it is a partial order and any $x,y\in S$ are comparable; then $(S,\leq_S)$ is an \hlt{ordered set}.
\item $C$ is a \hlt{chain} if $\leq_S|_C$ is an order.
\item $S$ is \hlt{well-ordered} if every nonempty $Y\subseteq S$ has $\ell\in Y$ with $\ell\leq_Sy$ for all $y\in Y$.
\item $\beta\in S$ is an \hlt{upper bound} of $E$ if $x\leq_S\beta$ for all $x\in E$; lower bounds reverse the inequality.
\item $E$ is \hlt{bounded above} if it has an upper bound; bounded below is dual.
\item $\sup E$ is the least upper bound of $E$, and $\inf E$ is the greatest lower bound of $E$, when these exist.
\item $m\in E$ is \hlt{maximal} if no $x\in E$ satisfies $m<_Sx$; minimal is dual.
\item $S$ has the \hlt{least-upper-bound property} if every nonempty bounded-above $E\subseteq S$ has $\sup E\in S$.
\end{enumerate}
\end{definition}

\begin{figure}[H]
  \centering
  \begin{tikzpicture}[
  x=0.80cm,
  y=0.82cm,
  >=Stealth,
  every node/.style={font=\small},
  axis/.style={gray!70, line width=0.55pt},
  setbar/.style={blue!70!black, line width=2.3pt, line cap=round},
  bound/.style={gray!50, line width=0.5pt},
  setpoint/.style={circle, fill=blue!70!black, inner sep=1.25pt},
  arrow/.style={->, gray!65, line width=0.5pt}
]
  \draw[axis] (-4.7,0) -- (4.7,0);

  \draw[setbar] (-1.35,0) -- (1.55,0);
  \foreach \x in {-0.75,-0.25,0.3,0.8,1.2} {
    \node[setpoint] at (\x,0) {};
  }

  \draw[arrow] (-1.35,0.60) -- (-1.35,0.16);
  \draw[arrow] (1.55,0.60) -- (1.55,0.16);
  \node[above] at (-1.35,0.60) {$\inf A$};
  \node[above] at (1.55,0.60) {$\sup A$};
  \node[below, yshift=-5pt] at (0.1,0) {$A$};

  \draw[bound] (-4.05,-0.42) -- (-4.05,-0.24) -- (-1.35,-0.24) -- (-1.35,-0.42);
  \node[below] at (-2.7,-0.42) {lower bounds};

  \draw[bound] (1.55,-0.42) -- (1.55,-0.24) -- (4.05,-0.24) -- (4.05,-0.42);
  \node[below] at (2.82,-0.42) {upper bounds};
\end{tikzpicture}

  \caption{The supremum is the least upper bound, while the infimum is the greatest lower bound.}
  \label{fig:supremum-infimum-bounds}
\end{figure}

\begin{example} \label{ex:supremum-need-not-be-attained} \hlt{Suprema Need Not Be Elements}\\
In $\Q$, the set $E_1=\{r\in\Q:r<0\}$ and the set $E_2=\{r\in\Q:r\leq0\}$ both have supremum $0$, but $0\notin E_1$ and $0\in E_2$. If $E=\{1/n:n=1,2,3,\ldots\}$, then $\sup E=1$ while $\inf E=0$, and the infimum is not a member of $E$.
\end{example}

\begin{theorem} \label{thm:lub-property-gives-glb-property} \hlt{Least Upper Bounds Give Greatest Lower Bounds}\\
If $S$ has the least-upper-bound property, then every nonempty bounded-below set $B\subseteq S$ has an infimum.
\begin{proof}
Let $L$ be the set of lower bounds of $B$. Then $L$ is nonempty and every $x\in B$ is an upper bound of $L$, so $\alpha=\sup L$ exists. Since every $x\in B$ bounds $L$ above, $\alpha\leq x$ for all $x\in B$; hence $\alpha$ is a lower bound of $B$.\\
If $\beta>\alpha$, then $\beta\notin L$, so $\beta$ is not a lower bound. Thus $\alpha=\inf B$.\\
\hlt{Structural Commentary}: Infima are suprema of lower bounds.\\
\hlt{Key Insight}: Completeness above gives completeness below.
\end{proof}
\end{theorem}

\begin{definition} \label{def:lattice-vocabulary} \hlt{Lattice Vocabulary}\\
Let $(L,\leq)$ be a partially ordered set.
\begin{enumerate}[label=(\roman*)]
\setlength{\itemsep}{0pt}
\item A \hlt{meet} of $x,y\in L$ is $\inf\{x,y\}$ when it exists, often written $x\wedge y$.
\item A \hlt{join} of $x,y\in L$ is $\sup\{x,y\}$ when it exists, often written $x\vee y$.
\item A \hlt{lattice} is a partially ordered set in which every two elements have a meet and a join.
\item A lattice is \hlt{complete} if every subset $E\subseteq L$ has an infimum and a supremum.
\item A \hlt{bottom} element is an element $\bot\in L$ with $\bot\leq x$ for every $x\in L$.\\
A \hlt{top} element is an element $\top\in L$ with $x\leq\top$ for every $x\in L$.
\item A lattice is \hlt{modular} if $x\leq z$ implies $z\wedge(y\vee x)=(z\wedge y)\vee x$ for all $x,y,z\in L$.
\end{enumerate}
\end{definition}
